%=====================================================================
\chapter{Managing Data Science Projects}\label{ch:projectmgmt}
%=====================================================================
\locator{6}{Part VI --- Professional Practice}

\begin{outcomes}
By the end of this chapter you will be able to:
\begin{tight}
  \item \textbf{Explain} why data science projects fail differently from software projects.
  \item \textbf{Structure} a project as a sequence of decision gates rather than a fixed plan.
  \item \textbf{Estimate} work whose outcome is genuinely unknown.
  \item \textbf{Select} metrics and build a dashboard that supports delivery decisions.
  \item \textbf{Identify} and \textbf{mitigate} the risks specific to AI-enabled projects.
\end{tight}
\end{outcomes}

\begin{prereq}
\chref{lifecycle} (CRISP-DM, framing, business versus model metrics),
\chref{mlops} (what ``done'' actually means) and \chref{ethics}.
\end{prereq}

\begin{keyterms}
project lifecycle $\bullet$ scope, time, cost, quality $\bullet$ agile $\bullet$
sprint $\bullet$ velocity $\bullet$ spike $\bullet$ decision gate $\bullet$
technical debt $\bullet$ reference class forecasting $\bullet$ leading and lagging
indicator $\bullet$ RAID log $\bullet$ stakeholder
\end{keyterms}

\section{Why Data Science Projects Are Different}

\begin{alertbox}[The Property That Breaks Normal Planning]
In a software project you know at the outset that the feature can be built; the
uncertainty is how long it takes. In a data science project you do not know whether
the thing is \emph{possible at all} until you have done a substantial part of the
work. The signal may not be in the data. No amount of planning discipline removes
this, and a plan that pretends otherwise will be wrong.
\end{alertbox}

\rowcolors{2}{white}{lightbg}
\begin{longtable}{@{}p{3.4cm}p{5.1cm}p{5.9cm}@{}}
\toprule
\rowcolor{primary!15}
& \textbf{Software project} & \textbf{Data science project} \\
\midrule
\endhead
Feasibility & Known at the start & \textbf{Unknown until you try} \\
Definition of done & Acceptance criteria pass & A metric threshold that may be unreachable \\
Estimation & Reference classes work reasonably & Highly unreliable before the data is understood \\
Main risk & Scope creep & \textbf{The data does not support the question} \\
After delivery & Maintenance & \textbf{Continuous retraining and monitoring} (\chref{mlops}) \\
Biggest cost & Development & Data acquisition and preparation (\dsref{effort-split}) \\
Failure mode & Late or over budget & Delivered, deployed, and quietly wrong \\
\bottomrule
\end{longtable}

\section{Gates, Not a Plan}

\dsfig{%
\begin{tikzpicture}[font=\scriptsize,
  ph/.style={rounded corners=3pt, draw=#1, fill=#1!8, text=#1!50!black,
             text width=2.5cm, align=center, font=\tiny\bfseries, minimum height=0.9cm},
  g/.style={diamond, draw=accent, fill=accent!10, text=accent!70!black, aspect=1.7,
            inner sep=1.5pt, font=\tiny\bfseries, align=center},
  arr/.style={-{Stealth[length=2mm]}, primary!55, line width=0.95pt},
  kill/.style={-{Stealth[length=2mm]}, accent!70, line width=0.9pt, dashed}]

\node[ph=secondary] (a) at (0,0)     {FRAME\\{\tiny 1 week}};
\node[g]            (g1) at (2.5,0)  {worth\\doing?};
\node[ph=tealhl]    (b) at (5.0,0)   {DATA\\FEASIBILITY\\{\tiny 2 weeks}};
\node[g]            (g2) at (7.6,0)  {signal\\present?};
\node[ph=goldhl]    (c) at (10.2,0)  {BASELINE\\MODEL\\{\tiny 2 weeks}};
\node[g]            (g3) at (12.8,0) {beats\\baseline?};

\node[ph=greenhl]   (d) at (2.5,-2.4)  {PRODUCTION\\MODEL\\{\tiny 4 weeks}};
\node[g]            (g4) at (5.4,-2.4) {ethics \&\\ops ready?};
\node[ph=purplehl]  (e) at (8.3,-2.4)  {DEPLOY\\{\tiny shadow $\to$ canary}};
\node[ph=primary]   (f) at (11.4,-2.4) {OPERATE\\{\tiny indefinitely}};

\foreach \x/\y in {a/g1, g1/b, b/g2, g2/c, c/g3}{\draw[arr] (\x) -- (\y);}
\draw[arr] (g3.south) .. controls +(0,-1.2) and +(0,1.2) .. (d.north);
\foreach \x/\y in {d/g4, g4/e, e/f}{\draw[arr] (\x) -- (\y);}

\node[ph=accent, minimum height=0.7cm] (stop) at (7.6,1.75) {STOP\\{\tiny and say so}};
\foreach \n in {g1,g2,g3}{\draw[kill] (\n) -- (stop);}

\node[anchor=north west, align=left, text width=13.6cm, font=\scriptsize, text=primary!60!black]
     at (-1.5,-3.85)
     {\textbf{The gates are the management technique.} Each is a real decision point with
      an explicit criterion agreed in advance, and each is allowed to end the project.
      Total spend before the first go/no-go is one week; before the second, three. A
      project killed at gate 2 for three weeks of effort is a \emph{success} of the
      process --- the failure is the one that runs nine months and then discovers the
      signal was never there.};
\end{tikzpicture}}%
{A data science project as a sequence of decision gates. Because feasibility is unknown at
the start, the plan must be structured to discover it cheaply and to permit an honourable
stop.}{project-gates}

\begin{conceptbox}[Make Stopping Respectable]
Teams continue doomed projects because stopping looks like failure. Fix this with
process: name the gates and their criteria in the project charter, so that stopping
at a gate is the plan working rather than the team losing. A group that has killed
projects at gate 2 is more trustworthy than one that has never killed anything.
\end{conceptbox}

\section{Estimation Under Genuine Uncertainty}

\begin{worked}[Estimating a Task Nobody Can Estimate]
\emph{``How long to build the dropout model?''}

\smallskip
\textbf{The wrong answer:} ``six weeks'' --- a single number, which will be treated
as a commitment and will be wrong.

\smallskip
\textbf{Better --- decompose by uncertainty type.}
\begin{center}
\begin{tabular}{@{}p{4.3cm}p{2.6cm}p{6.2cm}@{}}
\toprule
\textbf{Activity} & \textbf{Estimate} & \textbf{Why this confidence} \\
\midrule
Data access and permissions & 1--4 weeks & Depends on other people; historically the
longest pole \\
Data audit and cleaning & 2 weeks $\pm$1 & Similar to three past projects \\
Feature engineering & 1--3 weeks & Depends on what the audit reveals \\
Baseline model & 3 days & Genuinely routine \\
Reaching the target metric & \textbf{unbounded} & \textbf{May be impossible} \\
Deployment and monitoring & 2 weeks & Reference class exists \\
\bottomrule
\end{tabular}
\end{center}

\smallskip
\textbf{What to actually commit to.} Not the outcome, but the \emph{gate}:
\emph{``In three weeks we will report whether the data contains enough signal to
reach 70\% recall at 20\% flag rate. If it does, delivery is a further 6--8 weeks.
If it does not, we will recommend stopping and explain why.''}

\smallskip
\textbf{Why this is better management, not evasion.} It commits to a date, a
deliverable and a decision --- all of which are within the team's control --- and
declines to commit to a result that is not. A stakeholder given this can plan; a
stakeholder given ``six weeks'' cannot, because the number is fiction.
\end{worked}

\begin{definitionbox}[Spike]
A strictly time-boxed investigation whose deliverable is \emph{knowledge}, not
working software: ``spend three days determining whether the clickstream contains
week-4 signal.'' The time box is the point --- it converts open-ended research into
a bounded commitment.
\end{definitionbox}

\section{Metrics and Dashboards}

\dsfig{%
\begin{tikzpicture}[font=\scriptsize,
  hd/.style={rounded corners=3pt, fill=#1, text=white, font=\tiny\bfseries,
             text width=4.15cm, align=center, minimum height=0.5cm},
  b/.style={rounded corners=3pt, draw=#1, fill=#1!7, text width=4.15cm, align=left,
            font=\tiny, inner sep=5pt, minimum height=2.5cm}]

\node[hd=secondary] at (0,0)     {LEADING --- delivery health};
\node[b=secondary]  at (0,-1.5)  {%
  Available \emph{now}; predicts trouble.\\[3pt]
  $\bullet$ gate criteria met / missed\\
  $\bullet$ data access blockers open\\
  $\bullet$ cycle time per experiment\\
  $\bullet$ share of effort on rework\\[3pt]
  \emph{Act on these.}};

\node[hd=goldhl] at (4.75,0)     {LAGGING --- model quality};
\node[b=goldhl]  at (4.75,-1.5)  {%
  Available at each gate.\\[3pt]
  $\bullet$ PR-AUC vs baseline\\
  $\bullet$ recall at the agreed flag rate\\
  $\bullet$ subgroup performance (\chref{ethics})\\
  $\bullet$ stability across CV folds\\[3pt]
  \emph{Decide gates on these.}};

\node[hd=greenhl] at (9.5,0)     {OUTCOME --- did it work?};
\node[b=greenhl]  at (9.5,-1.5)  {%
  Available only after deployment.\\[3pt]
  $\bullet$ students retained\\
  $\bullet$ adviser hours per student helped\\
  $\bullet$ share of flags acted on\\
  $\bullet$ cost per retained student\\[3pt]
  \emph{Justify the work on these.}};

\node[anchor=north west, align=left, text width=13.8cm, font=\scriptsize, text=accent!75!black]
     at (-2.2,-3.15)
     {\textbf{``Share of flags acted on'' is the metric that predicts failure earliest.}
      A model whose output is ignored has already failed, however good its PR-AUC, and
      you can see that within a fortnight of launch --- long before any outcome data
      exists. Every deployed model should have one metric of this kind: not how good the
      predictions are, but whether anyone is using them.};
\end{tikzpicture}}%
{Three kinds of project metric, ordered by when each becomes available. Leading indicators
are for acting on, lagging for deciding gates, outcome measures for justifying the
work.}{project-metrics}

\section{Risk}

\rowcolors{2}{white}{defbg}
\begin{longtable}{@{}p{3.9cm}p{4.1cm}p{6.0cm}@{}}
\toprule
\rowcolor{primary!15}
\textbf{Risk} & \textbf{Likelihood} & \textbf{Mitigation} \\
\midrule
\endhead
The data lacks the signal & \textbf{High} & Gate 2 at three weeks; agree in advance that this ends the project \\
Data access takes months & \textbf{High} & Start the request on day one, in parallel with everything else \\
Leakage inflates results & High & Independent review of the split and features (\chref{wrangling}) \\
Stakeholder expects certainty & High & Communicate in ranges and probabilities from the first meeting (\chref{communication}) \\
The model is unfair to a subgroup & Moderate & Subgroup metrics in the gate criteria, not after deployment \\
Nobody uses the output & Moderate & Frame around the decision (\chref{lifecycle}); involve the user from gate 1 \\
Model decays after launch & \textbf{Certain} & Budget ongoing MLOps effort from the start (\chref{mlops}) \\
Key person leaves & Moderate & Reproducible pipelines, not notebooks (\chref{research}) \\
\bottomrule
\end{longtable}

\begin{alertbox}[Budget for After Launch]
The commonest structural error in data science project planning is treating
deployment as the end. A deployed model needs monitoring, periodic retraining,
incident response and eventual retirement. Plan roughly 20--30\% of the build
effort per year, indefinitely. A project plan that ends at deployment is planning a
system that will silently decay (\dsref{concept-drift}).
\end{alertbox}

%---------------------------------------------------------------------
\begin{fourlenses}{Managing the Work}
\lensPM{This chapter is the project-management lens in full. The distilled version:
gate the project so infeasibility is discovered in weeks rather than months, commit
to decisions rather than to outcomes, and budget for the operating phase from the
beginning.}

\lensLA{Education projects have an immovable constraint the others do not: the
\emph{academic calendar}. A model that misses the start of semester is useless for
a year, and you cannot compress a cohort. Plan backwards from week 1 of the
semester, and note that gate criteria involving student outcomes take a full term to
evaluate --- so the first genuine outcome evidence arrives after the second
deployment.}

\lensIOT{The critical path runs through \emph{hardware procurement}, not software.
Sensors have lead times, sites need physical access, and connectivity contracts take
weeks. Pilot on 20 devices before committing to 4{,}000: fleet failure modes are
invisible at bench scale, and the cost of discovering them post-deployment is
measured in site visits. Budget for the fact that firmware updates across a fleet
are themselves a project.}

\lensSEC{Success is invisible --- the business case rests on incidents that did not
occur --- so agree the metrics before starting: mean time to detect, mean time to
respond, coverage of monitored assets, alerts per analyst. Detection engineering is
continuous maintenance rather than a delivery, so the operating budget is not
optional. And note the tension with \chref{research}: publishing your method helps
attackers, which constrains how openly the project can be documented.}
\end{fourlenses}

%---------------------------------------------------------------------
\begin{lab}[A Project Charter That Fits on One Page]
\begin{lstlisting}[basicstyle=\ttfamily\tiny, frame=none, backgroundcolor=\color{codebg}]
PROJECT        Early identification of at-risk first-year students
SPONSOR        Head of Department          DATE  2026-09-01

DECISION       Which students receive an outreach call in week 5.
               (If no decision changes, there is no project.)
USERS          Three academic advisers. Involved from gate 1, not shown a demo at the end.

SUCCESS        >= 70% of eventual failures identified, flagging <= 20% of the cohort.
BASELINE       Current rule ("missed two labs") achieves ~45% recall at 15% flag rate.
               The model must beat THIS, not chance.

GATES
  G1  wk 1   Framing agreed, data access requested, decision confirmed  -> go/no-go
  G2  wk 3   Signal check: can ANY model reach 60% recall at 20% flags?  -> go/no-go
  G3  wk 5   Baseline beaten on walk-forward validation                  -> go/no-go
  G4  wk 9   Subgroup fairness within 5pp; monitoring plan signed off    -> go/no-go
             Stopping at any gate is a valid, expected outcome.

OUT OF SCOPE   Grading, disciplinary use, any automated action without a human.
ETHICS         Equal opportunity as the fairness criterion (approved by X, 2026-09-01).
               Counterfactual explanation shown to the adviser. Student may contest.

TOP RISKS      1. Signal absent            -> G2 kills the project at 3 weeks
               2. Data access slow          -> requested day 1; escalation path named
               3. Leakage                   -> split and features reviewed by Y at G3

AFTER LAUNCH   0.5 day/week monitoring; annual retrain per cohort; owner named: Z.
               Review for retirement after 3 years.
\end{lstlisting}

\textbf{Everything on this page is a commitment the team can keep.} Note what is
absent: a promised accuracy figure, and a delivery date for a result nobody can yet
know is achievable.
\end{lab}

%---------------------------------------------------------------------
\begin{checkpoint}
\begin{tightnum}
  \item Why can a data science project not be planned the way a software feature
        can? Name the specific property.
  \item A sponsor asks for a delivery date and an accuracy figure. What do you
        commit to instead, and how do you explain the difference?
  \item Your project is killed at gate 2 after three weeks. Was this a failure?
\end{tightnum}
\end{checkpoint}

%---------------------------------------------------------------------
\begin{chaptersummary}
\begin{tight}
  \item Data science projects differ because feasibility is unknown until you try;
        planning must be built around discovering that cheaply.
  \item Structure the work as decision gates with criteria agreed in advance, each
        able to end the project honourably.
  \item Estimate by decomposing into activities of different uncertainty, and commit
        to gates and deliverables rather than to results.
  \item Use time-boxed spikes for open-ended investigation.
  \item Track leading indicators to act, lagging indicators to decide gates, and
        outcome measures to justify the work.
  \item The largest planning error is treating deployment as the end. Budget 20--30\%
        of build effort per year for operations, indefinitely.
\end{tight}
\end{chaptersummary}

\begin{reviewq}
  \item \textbf{(MCQ)} The distinctive risk of a data science project is:
        (a)~scope creep (b)~the data may not support the question (c)~poor UI
        (d)~server capacity.
  \item \textbf{(MCQ)} A time-boxed investigation whose output is knowledge is a:
        (a)~sprint (b)~spike (c)~gate (d)~epic.
  \item \textbf{(Short)} Explain why a project killed at an early gate can be a
        success of the process.
  \item \textbf{(Short)} Give two leading and two lagging indicators for a data
        science project, and say how each is used differently.
  \item \textbf{(Short)} Why must the operating phase be budgeted at the start rather
        than requested later?
  \item \textbf{(Applied)} Write a one-page charter for a predictive maintenance
        project on 150 HVAC units, including four gates with explicit criteria, three
        risks with mitigations, and what you commit to versus what you decline to
        commit to.
  \item \textbf{(Applied)} A sponsor says ``the model must be 95\% accurate''.
        Explain, in terms they would accept, why that specification is both unhelpful
        and possibly meaningless, and propose a replacement.
\end{reviewq}
